\documentclass[11pt]{article}
\usepackage[a4paper,margin=2.2cm]{geometry}
\usepackage{xcolor}
\usepackage{fontspec}
\usepackage{xeCJK}
\usepackage{enumitem}
\setmainfont{Times New Roman}
\setCJKmainfont{SimSun}
\setlist[itemize]{leftmargin=1.4em,itemsep=0.35em,topsep=0.35em}
\setlength{\parindent}{0pt}
\setlength{\parskip}{0.55em}
\pagestyle{plain}
\begin{document}
{\color[HTML]{1F5FD2}\LARGE\bfseries 个别化教育计划（IEP）编写指南\par}
{\color[HTML]{667085}\small 星轨原创教育资源：用于家庭与学校共同制定支持计划。\par}
\vspace{0.8em}
{\color[HTML]{111827}\large\bfseries IEP 基本结构\par}
\begin{itemize}
\item 一份可执行的个别化教育计划应包含现有能力、长期目标、短期目标、服务安排、便利措施、评估方法、责任人和复盘日期。
\item 现有能力描述要来自自然情境中的可观察表现，不能只写诊断标签或一次测验分数。
\end{itemize}
{\color[HTML]{111827}\large\bfseries 可测量目标写法\par}
\begin{itemize}
\item 目标建议写清条件、行为和标准：在什么支持下，儿童完成什么行为，达到怎样的可测量水平。
\item 示例：给出视觉流程表后，儿童能在一周五天中至少四天独立完成三步入园流程。
\end{itemize}
{\color[HTML]{111827}\large\bfseries 进度记录\par}
\begin{itemize}
\item 每个目标只选择一种主要记录方式，例如次数、持续时间、等级量表、作品样本或观察记录。
\item 复盘时同时比较基线、当前表现和提示强度，决定是提高标准、调整策略还是延长练习周期。
\end{itemize}
\vfill
{\color[HTML]{667085}\footnotesize 本材料为应用内原创教育内容，不能替代医疗、法律或正式评估建议。\par}
\end{document}
